% Section 7: Conclusion

This paper formulates a mean-field game framework for modeling financial markets with heterogeneous overconfident investors. The model couples Bayesian belief updating (with misperceived signal precision) to endogenous price formation through a stochastic anchored-impact price equation. Our main pipeline is a myopic mean-field consistent equilibrium: we derive a checkable \CARA best response under Gaussian price increments and obtain a per-step stock--flow closure linking mean inventory/trading rates to mean belief. Separately, we record continuous-time \CARA/\LQG\ benchmarks as constructive comparisons.

The simulations demonstrate a clean \emph{mechanism separation}: (i) \emph{mispricing levels} are governed primarily by the microstructure regime $(\kappa,\lambda,\phi)$---weak anchoring and high noise trading generate large mispricing and volatility levels; and (ii) \emph{overconfidence ($k$)} robustly loads onto disagreement and trading intensity, but does not materially move price-level mispricing across a broad joint grid (Subsection~\ref{subsec:joint_grid}), even when mispricing levels become large. This structural separation---$k$ shifts belief dispersion and intensity, whereas $(\kappa,\lambda,\phi)$ govern mispricing magnitudes---is the main empirical takeaway. The identification argument (Section~\ref{sec:param_identification}) makes the model empirically testable, clarifying which observable moments pin down the structural parameters.

Several limitations and open questions remain. First, our baseline equilibrium concept and main simulation grid use the myopic \CARA best response under conditionally Gaussian price increments. As a validation check, Subsection~6.4.1 implements an unconstrained finite-horizon \CARA/\LQG\ intertemporal policy and finds similar qualitative $k$-effects. A fully constrained intertemporal formulation and a complete optimal-control MFG verification with endogenous filtering on prices remain future work. Second, while Subsection~\ref{sec:param_identification} provides an identification argument showing which moments pin down $(\kappa,\lambda,\sigma_v,\sigma_\epsilon,\phi)$ in principle, a full empirical calibration targeting these moments from price and volume data would strengthen the quantitative predictions. Third, the single-population framework abstracts from strategic and institutional heterogeneity; multi-population extensions (e.g., retail vs.\ institutional) could produce richer belief distributions and feedback. Finally, Section~III establishes well-posedness of the discrete-time particle system under the per-step contraction/uniqueness condition $0\le B_t<\Delta t$ (Theorems~\ref{thm:contraction} and~\ref{thm:well_posed}). A complete verification of existence/uniqueness conditions for the full intertemporal optimal-control MFG equilibrium (with endogenous filtering on prices) would strengthen the theoretical foundation.
